% Ghost-Desk Monitor: project proposal (Thai with English technical terms)
% Compile with XeLaTeX. On Overleaf: Settings (gear icon, bottom left) -> Compiler -> XeLaTeX.
% latexmkrc in this folder also swaps pdfLaTeX for XeLaTeX as a fallback.
% Font: Laksaman from the TLWG Thai fonts in TeX Live (a Sarabun-style face with Latin glyphs), so nothing needs uploading.
\documentclass[a4paper,11pt]{article}

\usepackage[a4paper,top=1.3cm,bottom=1.3cm,left=1.7cm,right=1.7cm]{geometry}
\usepackage{fontspec}
\usepackage{xcolor}
\usepackage{tabularx}
\usepackage{booktabs}
\usepackage{array}
\usepackage{enumitem}
\usepackage{titlesec}
\usepackage{tikz}
\usetikzlibrary{arrows.meta,positioning}
\usepackage[hidelinks]{hyperref}

% Thai text: break lines between Thai words (XeTeX uses ICU's Thai dictionary)
\XeTeXlinebreaklocale "th"
\XeTeXlinebreakskip = 0pt plus 1pt

% Script=Thai turns on the font's Thai mark positioning; without it vowels and tone marks sit too high or off-centre
\setmainfont{Laksaman}[
  Script = Thai,
  Language = Thai,
  Extension = .otf,
  UprightFont = *,
  BoldFont = *-Bold,
  ItalicFont = *-Italic,
  BoldItalicFont = *-BoldItalic,
  Scale = 1.13,
]
\setmonofont{TlwgTypewriter}[
  Extension = .otf,
  UprightFont = *,
  BoldFont = *-Bold,
  Scale = 1.05,
]

\linespread{0.95}
\setlength{\parindent}{0pt}
\setlength{\parskip}{3pt}
\renewcommand{\arraystretch}{1.0}
\setlength{\tabcolsep}{4pt}

\definecolor{muted}{gray}{0.35}
\titleformat{\section}{\large\bfseries}{\thesection.}{0.4em}{}
\titlespacing*{\section}{0pt}{5pt}{1pt}
\setlist{nosep,leftmargin=1.1em,itemsep=1pt}

% ragged-right columns: Thai has no spaces between words, so justified narrow cells stretch badly
\newcolumntype{L}{>{\raggedright\arraybackslash}X}
\newcolumntype{P}[1]{>{\raggedright\arraybackslash}p{#1}}
\newcolumntype{R}[1]{>{\raggedleft\arraybackslash}p{#1}}

\begin{document}

% ---------- title block ----------
{\LARGE\bfseries Ghost-Desk Monitor}\par
\vspace{1pt}
{\color{muted} ข้อเสนอโครงงาน (Project Proposal)}\par
\vspace{4pt}
{\small
\begin{tabular}{@{}>{\bfseries}l l@{}}
รายวิชา & Embedded System (CEDT) Final Project 2026: Smart Devices \\
กลุ่ม & TomyumSea \\
วันที่ & 30 ตุลาคม 2026 \\
Repository & \url{https://github.com/jiraphatnam2204/ghost-desk-monitor} \\
\end{tabular}}
\par\vspace{3pt}\hrule\vspace{2pt}

% ---------- 1 ----------
\section{ปัญหาและกลุ่มผู้ใช้ (Problem and Target Users)}

ในช่วงสอบ โต๊ะในห้องสมุดของมหาวิทยาลัยจำนวนมากถูกจองไว้ด้วยสิ่งของ (กระเป๋า หนังสือ) ขณะที่เจ้าของไม่อยู่เป็นเวลานาน
นิสิตที่ต้องการที่นั่งจึงหาที่นั่งไม่ได้ และเจ้าหน้าที่ห้องสมุดไม่มีทั้งข้อมูลและเวลาในการตรวจทุกโต๊ะ
ระบบจองที่นั่งด้วย QR ที่มีอยู่บันทึกได้เพียงว่าโต๊ะ\textbf{ถูกจอง} แต่ไม่รู้ว่ามี\textbf{ผู้นั่งอยู่จริง}หรือไม่
สมาชิกในทีมพบเหตุการณ์นี้เป็นประจำในห้องสมุดคณะช่วงสอบ

Ghost-Desk Monitor ติด sensor node ที่ทุกโต๊ะเพื่อตรวจว่ามีผู้นั่งอยู่จริงหรือไม่ และหมุนป้ายสามหน้า (ใช้งาน / ไม่อยู่ / ว่าง)
โดยไม่ต้องจอง ไม่ต้องเช็กอิน และไม่ต้องล็อกอินแอป โต๊ะที่ถูกทิ้งไว้พร้อมสิ่งของจะได้รับข้อความเตือนทาง LINE เมื่อครบ 20 นาที
และถูกปล่อยเมื่อครบ 30 นาที ส่วนโต๊ะที่เจ้าของเก็บของไปแล้วจะถูกปล่อยทันที

{\footnotesize
\begin{tabularx}{\linewidth}{@{}P{3.1cm}LL@{}}
\toprule
\textbf{กลุ่มผู้ใช้} & \textbf{ความต้องการ} & \textbf{สิ่งที่ระบบให้} \\
\midrule
นิสิตที่หาที่นั่ง & รู้ทันทีว่าโต๊ะใดว่างจริง & ป้ายสามหน้าบนทุกโต๊ะ และผังห้องแบบ real-time บนเว็บ \\
นิสิตที่ใช้โต๊ะอยู่ & ลุกไปชั่วครู่ได้โดยไม่เสียโต๊ะ & สถานะ ``ไม่อยู่'' และข้อความเตือนทาง LINE เมื่อครบ 20 นาที \\
เจ้าหน้าที่ห้องสมุด & รู้ว่าโต๊ะใดมีของค้าง ห้องแออัดหรือเสียงดัง โดยไม่ต้องเดินตรวจ & หน้า staff แสดงโต๊ะที่มีของค้าง ค่า CO\textsubscript{2} และระดับเสียง \\
\bottomrule
\end{tabularx}}

% ---------- 2 ----------
\section{แผนภาพบล็อก (Block Diagram)}

\begin{center}
\begin{tikzpicture}[
  font=\scriptsize,
  box/.style={draw, rounded corners=1pt, align=left, inner sep=3pt, text width=#1},
  arr/.style={-{Stealth[length=2mm]}, semithick},
  lab/.style={font=\tiny, align=center},
  every node/.append style={execute at begin node=\setlength{\baselineskip}{1.05em}},
]
\node[box=4.4cm] (node) {{\bfseries Sensor node $\times$ 2 (โต๊ะละตัว)}\\
  FireBeetle 2 ESP32-S3\\ LD2410 mmWave (UART, OUT)\\ VL53L1X ToF (I\textsuperscript{2}C)\\ OV2640 ถ่ายภาพเมื่อ gateway ขอ\\ Servo 270° $\rightarrow$ ป้ายปริซึม};
\node[box=4.1cm, right=2.0cm of node, line width=0.9pt] (gw) {{\bfseries Gateway $\times$ 1}\\
  YD-ESP32-23 (ESP32-S3)\\ SCD41 CO\textsubscript{2} (I\textsuperscript{2}C)\\ BH1750 แสง (I\textsuperscript{2}C)\\ MAX4466 ไมค์ (ADC1)\\ AI local + state machine};
\node[box=3.9cm, right=1.2cm of gw, yshift=6mm] (cloud) {{\bfseries Cloud}\\ Supabase: Postgres, REST, Realtime\\ Vision API (ภาพจาก gateway)};
\node[box=3.9cm, below=5mm of cloud] (dash) {{\bfseries Dashboard}\\ ผังห้อง + หน้า staff};
\node[box=4.1cm, below=4mm of gw] (line) {{\bfseries LINE Official Account}\\ broadcast เตือนนาทีที่ 20};

\draw[arr] ([yshift=3mm]node.east) -- node[lab, above]{ESP-NOW\\ ทุก 1 วิ} ([yshift=3mm]gw.west);
\draw[arr] ([yshift=-3mm]gw.west) -- node[lab, below]{SET\_SIGN,\\ SNAP, ack} ([yshift=-3mm]node.east);
\draw[arr] (gw.east |- cloud.west) -- node[lab, above]{HTTPS} (cloud.west);
\draw[arr] (cloud.south) -- node[lab, right]{Realtime} (dash.north);
\draw[arr] (gw.south) -- node[lab, right]{Messaging API} (line.north);
\end{tikzpicture}
\end{center}
\vspace{-4pt}

\begin{itemize}
\item \textbf{Sensor node} อ่านเรดาร์ (พลังงานแยกตามช่วงระยะ 0.75 m) และระยะจาก VL53L1X ถึงเบาะเก้าอี้ 10 ครั้งต่อวินาที
  ส่งชุดข้อมูล 210 ไบต์ (มี node ID) ทุก 1 วินาทีผ่าน ESP-NOW ถ่ายภาพเมื่อ gateway ขอ และหมุนป้ายด้วย servo
\item \textbf{Gateway} รันโมเดล AI (Edge Impulse) กับข้อมูล 5 วินาทีล่าสุดของแต่ละโต๊ะ (person / empty / nearby)
  ดูแล state machine และตัวจับเวลา ส่ง SET\_SIGN กลับไปที่ node และเป็นบอร์ดเดียวที่ต่ออินเทอร์เน็ต
\item \textbf{Cloud} Supabase เก็บข้อมูลดิบผ่าน HTTPS · cloud vision model ดูภาพหนึ่งภาพตอนต้องตัดสิน
  แล้วตอบ empty / person / items\_only · LINE Messaging API ส่งข้อความเตือน
\item \textbf{Dashboard} หน้าเว็บบน GitHub Pages แสดงผังห้องแบบ real-time และหน้า staff
  (โต๊ะที่มีของค้าง CO\textsubscript{2} ความสว่าง ระดับเสียง)
\end{itemize}

\textbf{กติกาของโต๊ะ} โต๊ะเป็น \emph{ใช้งาน} เมื่อโมเดลบน gateway เห็นผู้นั่งต่อเนื่อง 60 วินาที เมื่อไม่เห็นผู้นั่ง 2 นาที gateway ขอภาพจาก node
แล้วส่งให้ cloud vision model: \texttt{empty} ปล่อยโต๊ะทันที, \texttt{items\_only} เปลี่ยนเป็น \emph{ไม่อยู่} เตือนทาง LINE ที่ 20 นาที
และปล่อยที่ 30 นาทีพร้อมแจ้ง staff ว่ามีของค้าง, \texttt{person} คงสถานะใช้งาน หาก cloud ไม่ตอบ ใช้กติกา 30 นาทีตามปกติ
ภาพถูกถ่ายเฉพาะตอนตัดสินและไม่ถูกเก็บไว้

% ---------- 3 ----------
\section{เซนเซอร์และแอคชูเอเตอร์ (Bill of Materials)}

{\footnotesize
\begin{tabularx}{\linewidth}{@{}P{5.6cm}P{1.4cm}R{0.95cm}R{1.2cm}L@{}}
\toprule
\textbf{อุปกรณ์} & \textbf{บอร์ด} & \textbf{จำนวน} & \textbf{บาท} & \textbf{หมายเหตุ} \\
\midrule
FireBeetle 2 ESP32-S3 + OV2640 & Node & 2 & 2,370 & MCU และเซนเซอร์ 3: กล้อง \\
HLK-LD2410 mmWave 24 GHz & Node & 2 & 360 & เซนเซอร์ 1: ตรวจผู้นั่งแม้นั่งนิ่ง \\
VL53L1X time-of-flight & Node & 2 & 420 & เซนเซอร์ 2: ระยะถึงเบาะเก้าอี้ \\
TD-7120MG servo 0--270° & Node & 2 & 880 & แอคชูเอเตอร์: หมุนป้ายสามหน้า \\
YD-ESP32-23 (ESP32-S3) & Gateway & 1 & 0 & มีอยู่แล้ว (บอร์ดจาก Lab 10) \\
SCD41 CO\textsubscript{2} & Gateway & 1 & 820 & เซนเซอร์ 4: CO\textsubscript{2} (หมด: ใช้ MH-Z19C 895 แทน) \\
MAX4466 ไมโครโฟน & Gateway & 1 & 65 & เซนเซอร์ 5: ระดับเสียง \\
GY-302 BH1750 & Gateway & 1 & 50 & เซนเซอร์ 6: ความสว่าง \\
อแดปเตอร์ 5 V 2 A & Node & 2 & 260 & gateway ใช้ที่ชาร์จ USB-C ที่มีอยู่ \\
สายไฟ ตัวเก็บประจุ กล่อง งานพิมพ์ & -- & -- & \textasciitilde 400 & ประมาณการ \\
\midrule
\textbf{รวม} & & & \textbf{\textasciitilde 5,625} & ประมาณคนละ 1,125 บาท \\
\bottomrule
\end{tabularx}}

{\scriptsize\color{muted} ราคาร้านในไทย 28--30 ต.ค. 2026 ไม่รวมค่าส่ง · เซนเซอร์ 6 ชนิด (node 3 รวมกล้องตาม FAQ ข้อ 6, gateway 3) ผ่าน M1 แม้ไม่นับกล้อง}

% ---------- 4 ----------
\section{หัวข้อ Embedded Engineering Depth ที่เลือก (Section 5, II′)}

{\footnotesize
\begin{tabularx}{\linewidth}{@{}P{3.9cm}L@{}}
\toprule
\textbf{หัวข้อ} & \textbf{แนวทางการพัฒนา} \\
\midrule
\textbf{E1} RTOS multitasking & node มี 5 task (เรดาร์ ToF สื่อสาร กล้อง ป้าย) ส่งค่าผ่าน queue และใช้ mutex คุม I\textsuperscript{2}C bus \\
\textbf{E2} Interrupt-driven I/O & ขา OUT ของ LD2410 $\rightarrow$ ISR $\rightarrow$ semaphore ปลุก task เรดาร์ (ไม่ใช้ busy polling) \\
\textbf{E4} Reliable comm. & ack จาก send callback ของ ESP-NOW + retry 3 ครั้ง และ gateway พักข้อมูลใน RAM ตอนเน็ตหลุด (store-and-forward) \\
\textbf{E5} Security & HTTPS (TLS) ไป Supabase, LINE, vision API · key และ token อยู่ใน \texttt{secrets.h} ที่ไม่ commit \\
\textbf{E7} On-device AI & โมเดล Edge Impulse จากข้อมูลเรดาร์ + ToF รันบน gateway ESP32-S3 (FAQ ข้อ 4) ต้องแม่นกว่ากติกา threshold บนชุดทดสอบเดียวกัน · ประเมิน 13 พ.ย. \\
\bottomrule
\end{tabularx}}

% ---------- 5 ----------
\section{บทบาทของสมาชิก (Roles)}

{\footnotesize
\renewcommand{\arraystretch}{1.25}
% member cells: one person per line, "name (student ID)"; \newline starts the next person
% \addlinespace after each role keeps its names visibly grouped with it
\begin{tabularx}{\linewidth}{@{}P{7.2cm}L@{}}
\toprule
\textbf{บทบาท} & \textbf{สมาชิก} \\
\midrule
System Architect \& Integration & จิรภัทร นามวงศ์ \\ \addlinespace[3pt]
Embedded Developer (Sensor Node) & เขตโสภณ ครองยุติ\newline ชากิร บากา \\ \addlinespace[3pt]
Embedded Developer (Gateway \& Cloud) & เขตโสภณ ครองยุติ\newline ชากิร บากา \\ \addlinespace[3pt]
UI/UX \& Dashboard Developer & ธนกฤต สุนทรเวชพงษ์\newline ก้องภพ แสงแก้ว \\ \addlinespace[3pt]
Project Manager \& Test Lead & ก้องภพ แสงแก้ว \\
\bottomrule
\end{tabularx}}

\end{document}
